% =====================================================================
%  L13 backup slides: inside the operators (iterator model)
%  Kept out of the main flow; show them if students ask
%  "what does next() actually do inside a scan / join / sort?"
% =====================================================================
\sectionsubtitle{Lecture 13 \textbullet\ What \texttt{next()} does inside each operator}
\section{Backup: Operator Internals}

\subsection{Selection and heap scan}
\begin{frame}[fragile]{Backup: Selection and Heap Scan}
  \begin{columns}[T]
    \column{0.47\textwidth}
      {\small\textbf{Selection} $\sigma_p$ \textcolor{vinuGray}{(our $\sigma$ S.value $>$ 100)}}
      \begin{itemize}\footnotesize\setlength{\itemsep}{0.2ex}
        \item Streaming, ``on the fly'': a little work per row.
        \item State: the predicate $p$.
        \item \texttt{open()} / \texttt{close()}: just call the child's.
      \end{itemize}
\begin{lstlisting}[language=pseudo,basicstyle=\ttfamily\scriptsize\color{vinuNavy}]
next():
  loop:
    r = child.next()
    if r is EOF: return EOF
    if p(r): return r
\end{lstlisting}
    \column{0.50\textwidth}
      {\small\textbf{Heap Scan} \textcolor{vinuGray}{(our Scan R, Scan S)}}
      \begin{itemize}\footnotesize\setlength{\itemsep}{0.2ex}
        \item A leaf: no child, reads a \emph{heap file} (rows in no order).
        \item \texttt{open()}: open the file, cursor at page 1, slot 1.
        \item \texttt{next()}: return the row at the cursor, move to the next slot;
              at the end of a page, go to the next page; after the last page, \texttt{EOF}.
        \item \texttt{close()}: close the file.
      \end{itemize}
      \begin{center}
      \begin{tikzpicture}[pg/.style={draw=vinuNavy,fill=vinuNavy!6,minimum width=13mm,minimum height=11mm}]
        \foreach \i in {1,2,3}{
          \node[pg] (pg\i) at (\i*1.6,0) {};
          \node[font=\tiny,text=vinuGray] at ($(pg\i.south)+(0,-1.5mm)$) {page \i};
          \foreach \k in {1,2,3}{\draw[vinuNavy!40] ($(pg\i.north west)+(1mm,{-\k*2.7mm})$) -- ++(11mm,0);}}
        \draw[-{Stealth},vinuRed,line width=1.2pt] ($(pg2.north west)+(-2.5mm,-5.4mm)$) -- ++(2.3mm,0);
        \node[font=\tiny,text=vinuRed,anchor=south] at ($(pg2.north west)+(-1.5mm,0)$) {cursor};
      \end{tikzpicture}
      \end{center}
  \end{columns}
  \note{From the textbook (R\&G 14.1.1):     No index and no sort order $\Rightarrow$ no shortcut: scan every page and test every row. Cost = one I/O per page.}
\end{frame}

\subsection{Nested loops join and external sort}
\begin{frame}[fragile]{Backup: Nested Loops Join and Sort}
  \begin{columns}[T]
    \column{0.49\textwidth}
      {\small\textbf{Nested Loops Join} $R \bowtie_p S$}
      \begin{itemize}\footnotesize\setlength{\itemsep}{0.1ex}
        \item For each row of R (outer), go through all of S (inner).
        \item State: the current R row \texttt{r}.
        \item When S runs out: restart S, move to the next R row.
      \end{itemize}
\begin{lstlisting}[language=pseudo,basicstyle=\ttfamily\scriptsize\color{vinuNavy}]
next():
  loop:
    if r is EOF: return EOF
    s = right.next()
    if s is EOF:                  // S finished
      right.close(); right.open() // restart S
      r = left.next()             // next R row
      continue
    if p(r, s): return (r, s)
\end{lstlisting}
    \column{0.48\textwidth}
      {\small\textbf{Sort} (2-pass external sort)}
      \begin{itemize}\footnotesize\setlength{\itemsep}{0.2ex}
        \item \texttt{open()} does the heavy work, and it is \textbf{blocking}:
          \begin{itemize}\scriptsize
            \item \textbf{Pass 0}: read what fits in memory, sort it, write a \emph{sorted run} to disk. Repeat.
            \item Open every run, load its first page.
          \end{itemize}
        \item \texttt{next()} = \textbf{Pass 1 (merge)}: return the smallest front row of all runs.
        \item \texttt{close()}: delete the runs, close the child.
      \end{itemize}
      {\scriptsize\textcolor{vinuGray}{Analogy: sort 1,000 exam papers on a desk that holds 100: make sorted piles of 100, then keep taking the top paper that comes first.}}
  \end{columns}
  \note{From the textbook (R\&G 14.4.1), bigger chunks mean far less I/O:     {\scriptsize Reserves (1,000 pages) $\bowtie$ Sailors (500 pages):\;
    row at a time $\approx$ 50M I/Os (\textasciitilde140 h);\; page at a time 501,000 I/Os (\textasciitilde1.4 h);\;
    block of 100 pages 6,000 I/Os (\textasciitilde1 min). Same idea as the vectorised model.}}
\end{frame}
